\section{Recursos y presupuesto}

El presente proyecto requiere recursos humanos, materiales, logísticos, de servicios y tecnológicos para el desarrollo del modelo predictivo, la aplicación de apoyo a la decisión y la evaluación con productores de La Libertad. La planificación de estos recursos permite estimar costos, organizar el financiamiento y garantizar la ejecución del cronograma propuesto.

\subsection{Recursos}

Los recursos se clasifican según el clasificador de gastos institucional en cinco rubros: personal, bienes, viajes, servicios y tecnológicos. Cada rubro se detalla a continuación.

\subsubsection{Personal}

El equipo humano está conformado por los dos autores del proyecto de tesis y el asesor de investigación. Los autores dedican aproximadamente 20 horas semanales durante seis meses a las actividades de recolección de datos, modelado, desarrollo de software, redacción del informe y sustentación. El asesor participa en la orientación metodológica, revisión de avances y validación de instrumentos. Por tratarse de un proyecto de tesis de pregrado, el costo de personal se registra en S/ 0,00 al ser autofinanciado mediante dedicación académica.

\subsubsection{Bienes}

Comprende materiales de escritorio y campo necesarios para la documentación y aplicación de instrumentos: papel bond A4, carpetas, lapiceros, formatos impresos de encuesta, memorias USB para respaldo de datos y consumibles de impresión. Estos bienes facilitan la recolección de información en campo y la organización documental del proyecto.

\subsubsection{Viajes}

Se prevén dos a tres visitas a la región La Libertad (valles de Vir\'u, Chao y Moche) para trabajo de campo con productores de palta Hass, aplicación de cuestionarios y validación de requerimientos de la aplicaci\'on m\'ovil. Los viajes incluyen pasajes interprovinciales Trujillo--La Libertad y vi\'aticos b\'asicos (alimentaci\'on y movilidad local).

\subsubsection{Servicios}

Incluye acceso a internet de banda ancha, telefon\'ia m\'ovil, impresi\'on y fotocopiado de instrumentos, anillado del informe final, y eventual consultor\'ia especializada en despliegue de APIs. Estos servicios sostienen el desarrollo t\'ecnico y la elaboraci\'on del documento de tesis.

\subsubsection{Tecnol\'ogicos}

Se utilizan equipos de c\'omputo port\'atil de los propios autores para programaci\'on en Python (FastAPI), desarrollo frontend en React, entrenamiento de modelos SARIMA, LSTM y SVR, y an\'alisis de datos. El software empleado es de c\'odigo abierto (Python, bibliotecas cient\'ificas, entornos de desarrollo), por lo que no se contemplan licencias comerciales. El costo de amortizaci\'on de hardware se registra en S/ 0,00 al usar equipos existentes.

\subsection{Presupuesto}

La Tabla~\ref{tab:presupuesto_consolidado} presenta el presupuesto consolidado del proyecto, con costos unitarios y totales expresados en soles peruanos (S/). Los montos corresponden a estimaciones razonables para un proyecto de tesis de ingenier\'ia de sistemas con componente de campo en La Libertad.

\begin{table}[H]
\centering
\caption{Presupuesto consolidado del proyecto de tesis}
\label{tab:presupuesto_consolidado}
\small
\begin{tabularx}{\textwidth}{p{2.2cm}Xrrr}
\toprule
\textbf{Rubro} & \textbf{Descripci\'on} & \textbf{Cant.} & \textbf{C.U. (S/)} & \textbf{Total (S/)} \\
\midrule
Personal & Autores (2) y asesor --- dedicaci\'on acad\'emica & 6 meses & 0.00 & 0.00 \\
\midrule
Bienes & Papel bond A4 (resma) & 4 & 25.00 & 100.00 \\
Bienes & Carpetas y folders & 10 & 3.50 & 35.00 \\
Bienes & Formatos de encuesta impresos & 50 & 0.80 & 40.00 \\
Bienes & Memoria USB 32 GB & 2 & 35.00 & 70.00 \\
\midrule
Viajes & Pasaje Trujillo--La Libertad (ida y vuelta) & 6 & 45.00 & 270.00 \\
Viajes & Vi\'aticos y movilidad local & 3 & 80.00 & 240.00 \\
\midrule
Servicios & Internet y telefon\'ia m\'ovil & 6 meses & 60.00 & 360.00 \\
Servicios & Impresi\'on y fotocopiado & 1 & 180.00 & 180.00 \\
Servicios & Anillado y encuadernaci\'on informe & 3 & 25.00 & 75.00 \\
\midrule
Tecnol\'ogicos & Equipos de c\'omputo (uso propio) & 2 & 0.00 & 0.00 \\
Tecnol\'ogicos & Software libre (Python, React, etc.) & 1 & 0.00 & 0.00 \\
Tecnol\'ogicos & Servicio en la nube (despliegue prueba) & 6 meses & 50.00 & 300.00 \\
\midrule
\textbf{Total general} & & & & \textbf{2\,180.00} \\
\bottomrule
\end{tabularx}
\parbox{0.94\textwidth}{\small \textit{Nota.} Elaboraci\'on propia. C.U. = costo unitario.}
\end{table}

